\cotefile{53}
\scanpage{9}
% slanted
\margin{C'est \unsure{faux|sans|sous}} \deleted{\unsure{Coroll.|Coroll 1}} \uline{Corollaire} 2 \deleted{\unsure{Toute|Toutes}} \inserted{Les} \uline{composantes irréductibles de $G$} \deleted{(supposé loc. de t. f.)} \uline{sont les composantes connexes}. \uline{En particulier}, $G^{\unsure{0|\circ}}$ \uline{est irréductible}.

On peut supposer $k$ parfait, $G$ réduit donc normal (resp. lisse). Mais alors c'est \uline{vrai pour tt schéma normal} \inserted{(lisse)} sur $k$.... On a supposé ici $G$ loc. de t. f. sur $k$. Sinon on peut dire que $G_0$, la comp. irréd. de \unsure{$e$|$k$}, est stable par générisation, donc contient $\operatorname{Loc}_e(G)$.

% slanted
\margin{\textcircled{?} \deleted{C'est dans} le cas loc. de t. f. au moins}

% slanted
\margin{$-$ \unsure{$W^{\mathrm{no}}$|$W^{\mathrm{ro}}$|$N^{\mathrm{ro}}$} en $G$ \uline{loc. de t. f.}}

\deleted{\uline{Corollaire} 3 Pour \deleted{deux} ouverts $U$\inserted{, $V$} de $G$, \inserted{\deleted{loc. de t. f.} On peut s'en \inserted{dispenser} \deleted{supposer} $U$ \deleted{et} $V$ qu. cpct \uline{ou} $G$ loc. de t. f.)} $U\deleted{\unsure{U'|U}}V$ \deleted{contient $G^0$, si $U,U'$ sont des ouverts de $G$ qui contiennent $e$, $UV\supset G^0$.} \inserted{est une partie à la fois ouverte et fermée de $G$ (\unsure{évidemment|évidemt|évident} qu. cpct)}, donc stable par translation par $G^0$.}

\deleted{\deleted{\unsure{La deuxième assertion|La deuxième étant} conséquence immédiate de 2} première \deleted{\unsure{on le|on la} \unsure{pour|peut}} l'appliquant : $W=U\cap U^{-1}\cap V\cap V^{-1}$. \deleted{\illegible{2}} Pour prouver $UU^{-1}\supset G^0$, \deleted{\unsure{on peut}} \inserted{OPS par descente que $U$ contient pt rat $/k$, et par translation que $U\ni e$.} Soit $g\in G^0$, il faut montrer que (sur une ext $k'$ de \unsure{$k(g)$|$\kappa(g)$}) $g=\deleted{\unsure{g|g'}}uv^{-1}$, avec $u,v\in U(k')$, i.e. $(gU\cap U)(k')\ne \unsure{\phi|\varnothing}$, \deleted{et il suffit de prouver que (sur $k_1=k(g)$)} \deleted{$\unsure{gU\cap U\ne\varnothing|gU\cap U\ne e}$} or $gU\cap G^0$ et $U\cap G^0$ sont des ouverts non vides de $G^0$ (contient $g$ et $e$, donc ont une intersection non vide, puisque $G^0$ est irréd.}

\deleted{\uline{Autre dém} $G\times G \unsure{-|\longrightarrow} G$ est plat (car isom à $G\times G\xrightarrow{\mathrm{pr}_{1}}G$) de \deleted{\unsure{ses|sa}} générisant, donc l'image de $U\times U$ est stable par générisation, donc si $G$ est loc. de t. f.}

\deleted{La $=$ démonstration montre que $U.V$ stable par translation par $G^0$, donc étant qu. cpct il est \deleted{stable} ouvert et fermé.}

\deleted{$U.V$ ouvert comme image de $U\times_k V\subset G\times_k G$ par $G\times_k G \unsure{-|\longrightarrow} G$, qui est ouvert (isom : $G\times G\underset{\mathrm{pr}_1}{-}G$ qui l'est).}

\deleted{D'autre part, pour prouver que $U.V$ est fermé, il suffit, comme $U.V$ est \deleted{\unsure{vide ou quasi|vide ou qu. cpct}} qu. cpt donc rétrocompact (dans $G$ \uline{séparé}), \deleted{pour} de voir que $x\in U.V$, \deleted{$\unsure{\Rightarrow x\in U.V.|\Rightarrow y\in U.V.}$} \inserted{$y\in \unsure{(x)|\overline{(x)}|\overline{\{x\}}}$, alors $y\in U.V$.} On est ramené à \deleted{prouver} le prouver pour $x$}

% archive sheet number: 8
